\section{Introduction}
\label{sec:intro}

Industrial monitoring pipelines rarely stop at producing an anomaly score. Operators need a decision: when is a score high enough to trigger an inspection or a maintenance action? That decision requires a \emph{threshold}, which must be set without target-domain labels. The standard practice is to calibrate the threshold on source data and reuse it after deployment. We ask a simple question that, to our knowledge, has not been answered systematically: what happens to that threshold when the deployment domain differs from the training domain?

The question matters because cross-domain anomaly detection research has largely optimized the \emph{ranking} quality of detectors---autoencoders, contrastive domain adaptation, or time-series foundation models---while treating the decision threshold as an afterthought. Conformal prediction supplies distribution-free thresholds, but its exchangeability assumption is violated by both temporal dependence and domain shift; the adaptive-conformal literature has built machinery for temporal drift, yet a cross-\emph{domain} (cross-population) study on real industrial data is missing.

This paper reports such a study. We assemble six leave-one-group-out folds (dataset-level groups) over four domain families---semiconductor chemical-mechanical polishing, wind SCADA, photovoltaic plant measurements, and rotating machinery (three datasets)---keep each dataset at its native sampling rate, and fix documented labeling or injection protocols. Two guardrails are pre-committed: folds whose scores rank anomalies below chance are excluded from detection-quality claims, and the synthetic-only fold is flagged in every table. On every fold we train the same plain reconstruction encoder on source-normal windows and then deploy the score without target labels. We compare four thresholding mechanisms: a static source quantile, online adaptive conformal inference (ACI) \cite{aci2021}, a band-spectral divergence-shifted variant, and a conditional (Mondrian) variant that adapts separate quantiles across divergence regimes.

The results are unambiguous on the threshold side and honest about the detector side. Static thresholds fail in \emph{both} directions: on the wind fold the source threshold over-alarms so heavily that normal-window coverage collapses to $0.18$, while five folds under-alarm with coverage $0.989$--$0.999$ at a nominal $0.90$ level. Online ACI restores coverage to within $0.015$ of nominal on five folds across four nominal levels, and its advantage is not shared by rolling-quantile tracking or an extreme-value tail fit. Two additional findings cut against our own initial hypotheses and we report them as negative results: a divergence-shifted quantile adds nothing over plain ACI, and a band-spectral divergence statistic does not predict the transfer cost across folds. Finally, calibration and detection quality decouple---on two folds whose scores rank anomalies \emph{below} chance (a property of the datasets and labeling protocols, not of our encoder; feature baselines rank equally badly), ACI still controls coverage, but the alarms are useless for detection. We state plainly that calibration cannot rescue an anti-correlated ranker; coverage control is not a substitute for discriminative scores.

Our contributions are:
\begin{itemize}
\item \textbf{Benchmark.} A six-fold, four-domain-family leave-one-domain-out benchmark for cross-domain industrial anomaly detection at native sampling rates, with fully documented labeling and synthetic-injection protocols, provenance-checked datasets, and released code and configurations at \url{https://github.com/JibiaoZhu/crossdomain-conformal-calibration} (archived at \url{https://doi.org/10.5281/zenodo.23321391}).
\item \textbf{Calibration study.} Evidence that static source thresholds are unsafe under cross-domain shift in both directions, and that an online adaptive-conformal layer restores normal-window coverage on five of six folds, with bootstrap confidence intervals, an update-rate ablation, and calibration baselines (rolling quantile, extreme-value tail, slower updates).
\item \textbf{Method.} We identify an \emph{alarm-budget defect} of online ACI on anomaly-rich streams---the budget is consumed by anomalies and normal-window coverage over-shoots---and propose normal-anchored ACI (NA-ACI), which restricts the quantile feedback to likely-normal windows. NA-ACI reallocates alarms to anomalies (machinery-fold detection $0.12\!\to\!0.90$ with essentially unchanged FPR) and raises mean range-$F_1$ from $0.26$ to $0.39$ at comparable coverage, with an assumption-based FPR bound and an online exclusion-rate diagnostic; a single parameter interpolates toward ACI on anomaly-sparse folds.
\item \textbf{Negative results and diagnostics.} Divergence-shifted and conditional (Mondrian) threshold variants do not improve over ACI in this study; band-spectral divergence does not predict threshold inflation; and mechanism-heavy detector components (invariance losses, test-time adaptation, source gating) proved marginal in earlier ablations, which is why this paper studies calibration rather than architecture.
\item \textbf{Operational analysis.} False-alarm rates, detection delays, threshold traces, transfer-cost inflation ($0.5\times$--$19\times$), and per-run runtime/memory of the calibration layer ($0.6$--$0.8\,\mu$s per decision, 26\,MB peak), supporting a deployment guideline.
\end{itemize}

Section~\ref{sec:related} positions the work against cross-domain detectors and conformal prediction for time series. Section~\ref{sec:benchmark} defines the problem, metrics, datasets, and protocols. Section~\ref{sec:method} describes the calibration framework and its variants. Section~\ref{sec:experiments} presents the main results and analyses. Section~\ref{sec:discussion} distills the deployment guideline and limitations, and Section~\ref{sec:conclusion} concludes.
